\section{Conclusions}
\label{sec:conclusion}

As AI computing drives a rapid growth of data center electricity demand, AI data centers are increasingly expected to
respond to several DR events within one evening peak. This paper proposed a coordinated scheduling method for
consecutive DR in AI data center integrated energy systems that jointly recovers the computing backlog and the storage
across events, from the day-ahead commitment to intraday operation. Through theoretical analysis and case studies on a
16,000-GPU data center with 52 out-of-sample operating days, the main conclusions are summarised as follows.
\begin{enumerate}[label=(\arabic*),leftmargin=1.9em,itemsep=1.5pt]
\item Consecutive events leave a two-sided debt, deferred computing and depleted storage, that competes for one
connection-limited recovery channel. Independent recovery must repay the whole debt and retains only 38--39\% of the
next event's stand-alone capability for gaps of up to 1~h, whereas joint recovery repays only what the next event needs
and restores 96\% within 0.75~h from the same post-event state. The lag-bounded backlog constraint is exact
(Lemma~\ref{lem:edf}); the capability bound held in all 80 cases, the closed form of the recovery value matched the LP
within $10^{-5}$~MW, and the single channel price was confirmed with KKT residuals of $2\times10^{-12}$.
\item The main value of joint recovery is the reliably deliverable capacity. It raises the capacity that is delivered on
at least 95\% of the days from 6 to 8~MW per event (+33\%), which is equivalent to about 9--12~MWh of battery or 2.9~MW
of connection capacity, while price-coordinated independent recovery reaches this reliability at no commitment level.
The capability region of joint recovery is the single-event box intersected with one cross-event energy budget and can
be declared by four numbers; independent recovery shrinks this budget by 18\% even with a hindsight-optimal split, and a
one-shot price delivers only 60\% of the uniform capability on average because the IT side's response is not unique.
\item In closed-loop operation, joint recovery keeps all three events valid on all 52 days (benchmarks: 86--98\% of the
days) and lowers the total cost by 2.9--4.0\% against the strongest benchmark of each setting ($p<10^{-9}$), and by
14--16\% when the commitments approach the single-event capabilities; coupled commitment with joint recovery is 12.6\%
cheaper than single-event assessment with independent recovery. The saving comes from the rebound charge (78\%) when
rebound is priced and from the additional DR payment (86\%) under the Jiangsu settlement alone, not from energy
arbitrage.
\item Joint recovery can be distributed: when the IT scheduler and the facility controller exchange only power plans
and prices, 1, 3 and 10 exchanges per 15-min step recover 67\%, 85\% and 97\% of its advantage over the strongest
independent benchmark, and the converged price equals the recovery-channel price. Distributionally robust commitments
lower both the mean (by up to 0.26\%) and the worst-decile cost (by up to 0.4\%) out of sample, the scenario MPC adds
only an insignificant 0.4\% when the last event is uncertain, and the advantage of joint recovery holds across tariff
shapes ($-2.1$ to $-2.9\%$), seasons and climates ($-2.6$ to $-2.9\%$), liquid cooling ($-7.4\%$), thermal inertia
($-2.2\%$) and power capping, whereas tight IT ramp limits compress it (0.8\% at 1~MW per 15~min).
\item The fluid workload model is too optimistic for rigid gang-scheduled training: plans made with it fail at job level
on most days. With the rigid-job recovery bound the plans become executable (all events valid on 98\% and 86.5\% of the
days for the Acme and PAI loads), and joint recovery stays 3.0--5.0\% cheaper than the strongest benchmark at job level.
With behind-the-meter offshore wind the meter-based baseline reduces the deliverable commitments by 63\% at 20~MW, while
the advantage of joint recovery widens to 5.9--6.6\%; winter waste-heat recovery saves about 16\% of the daily cost
against boiler heating, and a hot-water tank turns the heat pump into a flexible load that is recharged after the events.
\end{enumerate}

Future work will extend the capability region to forecast errors and to the clearing of virtual-battery bids in DR
markets, schedule rigid jobs directly in the planning model and replace the first-order recovery bound by kernels based
on remaining run times, adopt Wasserstein ambiguity sets, and study the distributed coordination under communication
delays, asynchronous updates and strategic behaviour, as well as uncertainty in event timing and duration and
multi-site data center clusters with network constraints.
